\section{Resultado que se audita}
Sea $D$ el escalado bistocástico exacto de la máscara binaria basal de tamaño $j$.
Supondremos $\max jD_{rc}\le M_0$ y $\|D-J_j\|_{\op}\le r<1$, donde
$J_j=\mathbf1\mathbf1^t/j$. Para las máscaras de torres expresadas inicialmente
en su punto estacionario, el paso al escalado exacto conserva estos márgenes mediante
el rebalanceo del factor $1-P_e$ (R8; \emph{Media franjas}, parte 01).
Se quitan $v$ casillas marcadas para formar el exterior común. Para cada átomo se
quitan además $h$ filas y $h$ columnas; en la proyección local $h\le v$.
En la norma de activaciones global se usa $h\le aj$, con $a>0$ elegible y pequeño.
Cada menor real se escala por separado; llamamos $D'$ a ese escalado y $m=j-h$.
\paragraph{Teorema de los menores.}
Para parámetros $M_0,r$ fijos hay $a,\epsilon,c,C>0$ uniformes para los tableros
anteriores, siempre que $v/j$ sea suficientemente pequeño y $j$ suficientemente grande.
Para $Q=\lfloor\epsilon j\rfloor$ existe el aproximante explícito $H_Q$ de la parte 04, con
\[
 \left|\log\frac{m^m\per D'}{m!}-\log H_Q(1/m;D')\right|\le Ce^{-cj}.
\]
Para $1\le q\le Q$ la misma construcción de orden $q$ da
\[
 \left|\log\frac{m^m\per D'}{m!}-\log H_q(1/m;D')\right|
       \le (Cq/m)^q+Ce^{-cj}.
\]
Se usa el mismo índice de truncamiento en el padre y en cada cara. Los valores de
$1/m$ y de los núcleos balanceados sí corresponden al menor concreto.
En particular vale para $q=4v+12$, $v\le L\log n$, $j\ge\lceil\sqrt n\rceil$,
con $L$ fijo y $n$ bastante grande. En el disco $|u|\le c_0/q$, el logaritmo
del aproximante tiene rama holomorfa acotada para los datos complejos indicados
en la parte 02. La comparación con el conteo exacto se hace en los datos físicos reales.
\paragraph{Dependencias y alcance.}
Utilizamos la identidad del polinomio de torres centrado de Divoux--Wang v1,
secciones 3.1--3.3, la fórmula local de Wick de la prueba de su Proposición 3.8
y sus dos estimaciones en módulo, identificadas en FUENTES.md. Los órdenes
crecientes se deducen aquí. Para el logaritmo del error puro de torres se conserva
la normalización exacta de R9, partes 06--07. Las entradas geométricas de llegada
y las cancelaciones del predictor de reinas son cuestiones distintas.
Todas las afirmaciones de uniformidad requieren el mismo $M_0,r$ para la familia.
No se establece aquí que la ley de llegada visite exclusivamente esta familia.
